\partchapter{4}{第四部分\quad 总结}

\Needspace{85pt}

\section{可扩展之处}

以下方向均未实施，属后续计划，不计入本轮成果。

\begin{enumerate}
\def\labelenumi{\arabic{enumi}.}
\tightlist
\item
  长消息支持：当前 PT/AD 各 $\leq$16 B，经 AXI-Lite 字节 FIFO 传输；可改为共享缓冲或流式搬运，并补齐完整 1089 组板级 KAT 与多块吞吐测量。
\item
  器件绑定密钥：将独立 RO-PUF/Device DNA 原型并入当前 overlay，关闭派生输入回读通路，补充跨上电、温压与多器件统计后再评估密钥用途。
\item
  指令与算法：在现有 custom-0 状态指令上扩展 SP 800-232 中的 Ascon-Hash256/XOF128，并以 intrinsic 封装降低软件调用开销。
\item
  工程验证：补独立频率测量、小时级稳定性、完整 ISA/异常矩阵与干净环境复现。
\end{enumerate}

\Needspace{85pt}

\section{心得体会}

实现可运行与结论可复核是不同目标。标准 KAT、真实 CPU 仿真与实板测试各自覆盖不同问题，不能用某一层通过替代另一层验收；相同文件名也不能证明工件属于同一构建。因此，发布身份、逐文件哈希、source\_dirty 和源集选择应与原始样本一并保存。

性能比较首先是测量边界的比较。当前固定 16 B 输入下 MMIO 与 ISE 均降低了固件调用周期，但 PS 侧调用仍包含 Python 与通信开销，响应时间没有获得同倍数改善。保留这一区别，比仅展示最大的周期加速比更有解释力。

验证框架同样需要负向用例。错误长度、过期 epoch、重复请求及 FIFO overflow 的拒绝与恢复，为演示链的健壮性提供了有限但可复查的证据；它们不构成无限时稳定性、完整 ISA 合规或安全认证证明。报告应保留已知失败与未测项，不以文字更新代替工程验证。
